\subsection{A small middle corner opposite a binary stalk}

\begin{lemma}[Proved: a fixed middle-to-last arrow with trivial first corner]
\label{CP2:A-linear-joint-fixed-middle-arrow}
Let $C$ have full linear support $0<1<2$, trivial first
corner, finite middle and last corners $B,N$, and nonempty
forward hom-sets $X,Y,Z$.  If some $y\in Y$ is fixed by
both its right $B$ action and its left $N$ action, every
rational degree-two product in $BC$ vanishes.
\end{lemma}
\begin{proof}
Put $D=C|_{\{0,1\}}$ and $I=\{yx:x\in X\}\subset Z$.
For $z\in I$ put $X_z=\{x:yx=z\}$; this is a nonempty
left $B$-subact of $X$.  Define a category $D'$ with the
middle object of $D$, together with a separate trivial
initial object $0_z$ for each $z\in I$, and with hom-set
$X_z$ from $0_z$ to the middle object.  Its actual fold
$D'\to D$ identifies the objects $0_z$ and is the identity
on the union $X=\coprod_zX_z$ and on $B$.

On nerves this only identifies the $|I|$ distinct initial
vertices.  No nondegenerate simplex contains two such
vertices.  Hence it is the pushout identifying this finite
set of vertices to one, and the relative homology sequence
gives an isomorphism $H_j(BD';\mathbb Q)\to H_j(BD;\mathbb Q)$
for $j\geq2$.  Both spaces are connected.

Let $E$ be the category of elements of the connecting
right $D$-act.  The objects $(1,y)$ and $(0,z)$ for $z\in I$
form a full subcategory isomorphic to $D'$: its arrows
$(0,z)\to(1,y)$ are exactly the elements of $X_z$.
The left $N$ action fixes this subcategory pointwise.
Indeed it fixes $y$, it fixes $yx$, and it leaves the
$D$-arrow labels unchanged.  Thus the two-sided bridge
contains the actual bar map $BN\times BD'\to E_{hN}$.
Its composite to the two vertex blocks is the identity
on $BN$ and the fold on $BD'$.  The joint homology map
\[
 H_2(E_{hN};\mathbb Q)\longrightarrow
 H_2(BN;\mathbb Q)\oplus H_2(BD;\mathbb Q)
\]
is surjective.  Dually, both restrictions of every
degree-two collage class vanish by Mayer--Vietoris.
The two-block relative suspension argument annihilates
their products.
\end{proof}

\begin{theorem}[Proved: a small middle corner and two short incoming stalks]
\label{CP2:A-linear-four-middle-binary-short-long}
Let $C$ have full linear support with first corner of
order at most two, middle corner of order at most four,
and last corner of order at most eleven.  Suppose
$|C(1,2)|=2$ and $1\leq|C(0,2)|\leq3$.
If its rational homology is bounded, every rational
degree-two product vanishes.  There is no restriction
on the nonempty hom-set $C(0,1)$.
\end{theorem}
\begin{proof}
The actual exact-orbit quotient at the first corner is
a rational homology equivalence, since a monoid algebra
of order at most two is semisimple.  It preserves the
ring and boundedness, leaves $Y=C(1,2)$ unchanged, and
does not increase $|Z|$.  We may therefore suppose the
first corner trivial.  Write $B,N$ for the other corners.

Consider the transformation image of the right $B$ action
on the two-set $Y$.  If it is trivial, apply
Lemma~\ref{CP2:A-linear-trivial-middle-action-H2}.
Here its quotient functor has an actual section: choose
one representative in each $B$ orbit of $X$ and include
the identity of $B$, with $Y,Z,N$ unchanged.  All products
are respected because $\mu(y,x)$ is constant on each
$B$ orbit.  Thus the quotient has bounded homology.
Its first two corners are trivial, and
Proposition~\ref{CP2:A-two-small-incoming-eleven} kills
its second products.  Pull back along the $H^2$ isomorphism.

If exactly one constant transformation occurs, its value
is fixed by all right $B$ actors.  Every commuting left
$N$ actor fixes that value also.  The preceding fixed-arrow
lemma applies.

If the image is the permutation group $C_2$, the minimal
group of $B$ maps onto $C_2$.  The order bound $|B|\leq4$
forces its minimal ideal to have one row and one column:
in the nongroup case $rc|G|+1\leq4$, with $|G|\geq2$;
the group case is immediate.  Both trivial rational
$B$ modules are projective.  Exact orbits identify the
source block $D=C|_{\{0,1\}}$ rationally with a graph,
so $H^2(BD;\mathbb Q)=0$.  The commuting $N$ action on
$Y$ consists of permutations.  Apply
Lemma~\ref{CP2:A-permutation-fiber-cups}.

Finally suppose both constants occur, with or without
the interchange.  Every commuting left $N$ actor is
the identity on $Y$.  The minimal ideal of $B$ maps
onto the two constants, which form a right-zero band
when regarded as a quotient of the right-acting monoid.
Thus it has at least two columns.  Since $|B|\leq4$,
the inequality $rc|G|+1\leq4$ forces $r=|G|=1$.
The trivial left $\mathbb QB$ module is projective.

Let $E=C|_{\{1,2\}}$.  The $N$ action on its connecting
biset $Y$ is trivial.  Lemma~\ref{CP2:A-trivial-cut-H2}
identifies $H^2(BE)$ with that of the ordinary right
$B$-act cone on $Y$.  Its reduced homology in degrees
at least two is
$\operatorname{Tor}_{j-1}^{\mathbb QB}
(\ker(\mathbb QY\to\mathbb Q),\mathbb Q)$, which is
zero by the stated left projectivity.  Hence $H^2(BE)=0$.
The first corner is trivial, so both vertex blocks of
the cut $\{0\}\to E$ have zero second cohomology.
The relative suspension argument finishes this case.
These cases exhaust the binary transformation images.
\end{proof}

\begin{corollary}[Proved: two further large linear nineteen-arrow profiles]
\label{CP2:A-linear19-two-short-binary-large-profiles}
The profiles $(1,3,8;2,2,3)$ and $(1,3,8;3,2,2)$,
original indices $77,78$ in the frozen linear inventory,
cannot have classifying space of projective-plane type.
\end{corollary}
\begin{proof}
They satisfy the theorem, and projective-plane homology
is bounded with a nonzero rational degree-two square.
\end{proof}
